\chapter{Fixed-Cloud Filtering, Same-Branch SGQF, and Lane-Specific Validation}
\label{ch:bf-highdim-fixed-cloud-sgqf-validation}
\providecommand{\widehatell}{\widehat{\ell}}
\providecommand{\dotwidehatell}{\dot{\widehat{\ell}}}
\providecommand{\dd}{\,\mathrm d}

\section{Purpose And Scope}
\label{sec:bf-hd-sgqf-validation-purpose}

The previous chapter fixed the sparse-grid object itself: the one-dimensional
rule family, the active band, the signed Smolyak combination, and the merged
cloud stored in standardized Gaussian coordinates.  That was the construction
burden.  The present chapter takes up the burden that follows once the cloud has
already been fixed: what scalar the fixed cloud computes, how the forward value
path is assembled, what counts as the same branch for this lane, and what
validation burden must be discharged before the sparse-grid surrogate can be
promoted beyond a pedagogical or exploratory object.

This split is deliberate.  A reader first needs to understand how the cloud is
built.  Only after that does it become useful to ask which one-step scalar the
cloud defines, which branch metadata must stay fixed when differentiating that
scalar, and which diagnostics certify that a BayesFilter implementation is still
evaluating the same declared object rather than a nearby adaptive variant.

\section{Filtering Value Path On A Fixed Cloud}
\label{sec:bf-hd-sgqf-value-path}

Fix a standardized sparse-grid cloud
\begin{equation}
  \mathcal C=\{(\xi^{(r)},w_r)\}_{r=1}^M,
  \qquad
  \xi^{(r)}\in\R^{n_x},
  \quad
  w_r\in\R.
  \label{eq:bf-hd-sgqf-cloud-for-filter}
\end{equation}
and freeze that cloud before likelihood evaluation.  The sparse-grid value path
then has two quadrature stages at each time: prediction and observation.

\subsection{Prediction Stage}
\label{sec:bf-hd-sgqf-prediction-stage}

Let \(P_{t-1}=C_{t-1}C_{t-1}^\top\) on the declared square-root branch.  Place
the previous Gaussian cloud by
\begin{equation}
  \chi_{t-1}^{(r)} = m_{t-1}+C_{t-1}\xi^{(r)},
  \label{eq:bf-hd-sgqf-pred-placed-points}
\end{equation}
push points through the transition,
\begin{equation}
  a_t^{(r)}=f_\theta(\chi_{t-1}^{(r)}),
  \label{eq:bf-hd-sgqf-transition-points}
\end{equation}
and form
\begin{align}
  m_t^-&=\sum_{r=1}^Mw_ra_t^{(r)},
  \label{eq:bf-hd-sgqf-pred-mean}\\
  P_t^-&=Q_\theta+\sum_{r=1}^Mw_r
  (a_t^{(r)}-m_t^-)(a_t^{(r)}-m_t^-)^\top.
  \label{eq:bf-hd-sgqf-pred-cov}
\end{align}
After \eqref{eq:bf-hd-sgqf-pred-cov}, BayesFilter symmetrizes \(P_t^-\) and
checks the declared positive-definite branch.  For the same-branch lane studied
here there is no adaptive jitter, rank switch, node rewrite, or cloud rebuild
inside the scalar.  If the predictive branch fails, the value path returns a
veto instead of silently changing the object being evaluated.

\subsection{Observation Stage}
\label{sec:bf-hd-sgqf-observation-stage}

Factor \(P_t^-=C_t^-(C_t^-)^\top\) on the declared branch and place the
predictive cloud:
\begin{equation}
  \chi_t^{(r)} = m_t^-+C_t^-\xi^{(r)},
  \label{eq:bf-hd-sgqf-obs-placed-points}
\end{equation}
Propagate through the observation map,
\begin{equation}
  z_t^{(r)}=h_\theta(\chi_t^{(r)}),
  \label{eq:bf-hd-sgqf-obs-values}
\end{equation}
and compute
\begin{align}
  \bar z_t&=\sum_{r=1}^Mw_rz_t^{(r)},
  \label{eq:bf-hd-sgqf-obs-mean}\\
  \Delta_t^{(r)}&=z_t^{(r)}-\bar z_t,
  \label{eq:bf-hd-sgqf-obs-centered}\\
  S_t&=R_\theta+\sum_{r=1}^Mw_r
  \Delta_t^{(r)}\Delta_t^{(r)\top},
  \label{eq:bf-hd-sgqf-obs-cov}\\
  C_{xz,t}&=\sum_{r=1}^Mw_r
  (\chi_t^{(r)}-m_t^-)\Delta_t^{(r)\top}.
  \label{eq:bf-hd-sgqf-cross-cov}
\end{align}
The one-step innovation scalar is then
\begin{align}
  v_t&=y_t-\bar z_t,
  \label{eq:bf-hd-sgqf-innovation}\\
  \widehat\ell_t^{\rm FSGQ}
  &=-\frac12
  \left\{\log\det S_t+v_t^\top S_t^{-1}v_t+n_y\log(2\pi)\right\}.
  \label{eq:bf-hd-sgqf-loglik}
\end{align}
If \(S_t\) fails the declared branch, the value path returns a veto.  If it
passes, the Gaussian projection update is
\begin{align}
  K_t&=C_{xz,t}S_t^{-1},
  \label{eq:bf-hd-sgqf-kalman-gain}\\
  m_t&=m_t^-+K_tv_t,
  \label{eq:bf-hd-sgqf-filter-mean}\\
  P_t&=P_t^- - K_tS_tK_t^\top.
  \label{eq:bf-hd-sgqf-filter-cov}
\end{align}
with the final carried covariance again checked on the declared branch.

\paragraph{Value-path execution order.}
For this lane the runtime order is part of the scalar identity:
\begin{enumerate}[leftmargin=0.28in]
  \item symmetrize and factor \(P_{t-1}\),
  \item place previous-state points and propagate through \(f_\theta\),
  \item form \((m_t^-,P_t^-)\) and veto if the predictive branch fails,
  \item factor \(P_t^-\), place predictive points, and propagate through
  \(h_\theta\),
  \item form \(\bar z_t\), \(S_t\), and \(C_{xz,t}\), and veto if the
  innovation branch fails,
  \item form \(v_t\), \(\widehat\ell_t^{\rm FSGQ}\), \(K_t\), \(m_t\), and
  \(P_t\), then veto if the carried covariance branch fails.
\end{enumerate}
Any implementation that changes the cloud, node ordering, factor branch, or
veto policy in the middle of this sequence is no longer evaluating the same
fixed-cloud scalar.

\section{Worked One-Step Example With A Numeric Oracle}
\label{sec:bf-hd-sgqf-worked-example}

Use the one-dimensional transition-observation model
\begin{align}
  X_0 &\sim \calN(m_0,P_0),
  &m_0&=0.5,
  &P_0&=1,
  \label{eq:bf-hd-sgqf-example-initial}\\
  X_1 &= f_\beta(X_0)+\eta_1,
  &f_\beta(x)&=\beta x,
  &\eta_1&\sim \calN(0,Q),
  \label{eq:bf-hd-sgqf-example-transition}\\
  Y_1 &= h_\theta(X_1)+\varepsilon_1,
  &h_\theta(x)&=x^2,
  &\varepsilon_1&\sim \calN(0,R),
  \label{eq:bf-hd-sgqf-example-observation}
\end{align}
with
\begin{equation}
  Q=0.25,
  \qquad
  R=1,
  \qquad
  y_1=2.
  \label{eq:bf-hd-sgqf-example-QRy}
\end{equation}
Evaluate at \(\beta=1\) on the positive scalar Cholesky branch and use the
three-point standard-normal cloud
\begin{equation}
  \mathcal C_{1,2}
  =
  \left\{
  \left(-\sqrt3,\frac16\right),
  \left(0,\frac23\right),
  \left(\sqrt3,\frac16\right)
  \right\}.
  \label{eq:bf-hd-sgqf-example-cloud}
\end{equation}

\paragraph{Placed points and moments.}
Since \(C_0=\sqrt{P_0}=1\), the previous-state placed points are
\begin{equation}
  \chi_0^{(r)}=m_0+C_0\xi^{(r)}
  \in
  \{0.5-\sqrt3,\ 0.5,\ 0.5+\sqrt3\}.
  \label{eq:bf-hd-sgqf-example-prev-points}
\end{equation}
At \(\beta=1\), the transition images satisfy
\begin{equation}
  a_1^{(r)}=f_\beta(\chi_0^{(r)})=\chi_0^{(r)}.
  \label{eq:bf-hd-sgqf-example-transition-images}
\end{equation}
The predicted Gaussian moments are therefore
\begin{equation}
  m_1^-=
  \sum_rw_ra_1^{(r)}=0.5,
  \qquad
  P_1^-=Q+\sum_rw_r(a_1^{(r)}-m_1^-)^2=1.25=\frac54.
  \label{eq:bf-hd-sgqf-example-pred-moments}
\end{equation}
The declared Cholesky factor is
\begin{equation}
  C_1^- = \sqrt{P_1^-}=\sqrt{1.25}=\frac{\sqrt5}{2}.
  \label{eq:bf-hd-sgqf-example-pred-factor}
\end{equation}
Hence the predictive placed points are
\begin{equation}
  \chi_1^{(r)}=m_1^-+C_1^-\xi^{(r)}
  \in
  \left\{0.5-\frac{\sqrt{15}}{2},\ 0.5,\ 0.5+\frac{\sqrt{15}}{2}\right\}.
  \label{eq:bf-hd-sgqf-example-pred-points}
\end{equation}
Passing those points through \(h_\theta(x)=x^2\) gives
\begin{equation}
  z_1^{(r)}
  \in
  \left\{2.0635083269,\ 0.25,\ 5.9364916731\right\}.
  \label{eq:bf-hd-sgqf-example-z-values}
\end{equation}
The centered observation values are
\begin{equation}
  \Delta_1^{(r)}=z_1^{(r)}-\bar z_1
  \in
  \{0.5635083269,\ -1.25,\ 4.4364916731\}.
  \label{eq:bf-hd-sgqf-example-Delta}
\end{equation}
and the Gaussian update objects are
\begin{align}
  \bar z_1&=\sum_rw_r z_1^{(r)}=1.5=\frac32,
  \label{eq:bf-hd-sgqf-example-zbar}\\
  S_1&=R+\sum_rw_r(z_1^{(r)}-\bar z_1)^2=5.375=\frac{43}{8},
  \label{eq:bf-hd-sgqf-example-S}\\
  C_{xz,1}&=\sum_rw_r(\chi_1^{(r)}-m_1^-)(z_1^{(r)}-\bar z_1)=1.25=\frac54,
  \label{eq:bf-hd-sgqf-example-Cxz}\\
  v_1&=y_1-\bar z_1=0.5=\frac12,
  \label{eq:bf-hd-sgqf-example-v}\\
  K_1&=C_{xz,1}S_1^{-1}=\frac{10}{43}\approx0.2325581395,
  \label{eq:bf-hd-sgqf-example-K}\\
  m_1&=m_1^-+K_1v_1=\frac{53}{86}\approx0.6162790698,
  \label{eq:bf-hd-sgqf-example-m}\\
  P_1&=P_1^- -K_1S_1K_1^\top=\frac{165}{172}\approx0.9593023256.
  \label{eq:bf-hd-sgqf-example-P}
\end{align}
The one-step deterministic approximate log likelihood is
\begin{equation}
  \widehat\ell_1^{\rm FSGQ}
  =
  -\frac12\left\{\log S_1+v_1^2S_1^{-1}+\log(2\pi)\right\}
  \approx -1.7830736342.
  \label{eq:bf-hd-sgqf-example-ell}
\end{equation}

\paragraph{Same-branch derivative with respect to \texorpdfstring{\(\beta\)}{beta}.}
The same cloud and the same positive scalar Cholesky branch are reused in the
derivative lane.  Because the parameter enters only through the transition in
this example, \(\partial_\beta h_\theta(x)=0\).  Since \(m_0\) and \(P_0\) are
fixed in \(\beta\), the previous-state placed points satisfy
\begin{equation}
  \dot\chi_0^{(r)}=0.
  \label{eq:bf-hd-sgqf-example-prev-dot}
\end{equation}
From the transition derivative rule,
\begin{equation}
  \dot a_1^{(r)}
  =
  D_xf_\beta(\chi_0^{(r)})\dot\chi_0^{(r)}+\partial_\beta f_\beta(\chi_0^{(r)})
  =
  \chi_0^{(r)}
  \in
  \{0.5-\sqrt3,\ 0.5,\ 0.5+\sqrt3\}.
  \label{eq:bf-hd-sgqf-example-adot}
\end{equation}
Hence
\begin{equation}
  \dot m_1^-=
  \sum_rw_r\dot a_1^{(r)}=0.5,
  \qquad
  \dot d_1^{(r)}=\dot a_1^{(r)}-\dot m_1^-
  \in
  \{-\sqrt3,\ 0,\ \sqrt3\},
  \label{eq:bf-hd-sgqf-example-mdot-dotd}
\end{equation}
and
\begin{equation}
  \dot P_1^-=2.
  \label{eq:bf-hd-sgqf-example-Pminus-dot}
\end{equation}
Because \(C_1^-=(P_1^-)^{1/2}\), the declared branch derivative is
\begin{equation}
  \dot C_1^-=
  \frac{\dot P_1^-}{2C_1^-}
  =
  \frac{2}{\sqrt5}
  \approx 0.8944271910.
  \label{eq:bf-hd-sgqf-example-Cdot}
\end{equation}
Therefore the predictive point sensitivities are
\begin{equation}
  \dot\chi_1^{(r)}=\dot m_1^-+\dot C_1^-\xi^{(r)}
  \in
  \{-1.0491933385,\ 0.5,\ 2.0491933385\}.
  \label{eq:bf-hd-sgqf-example-chidot}
\end{equation}
Since \(h_\theta(x)=x^2\),
\begin{equation}
  \dot z_1^{(r)}=2\chi_1^{(r)}\dot\chi_1^{(r)}
  \in
  \{3.0143149884,\ 0.5,\ 9.9856850116\}.
  \label{eq:bf-hd-sgqf-example-zdot}
\end{equation}
The centered observation sensitivities are
\begin{equation}
  \dot\Delta_1^{(r)}=\dot z_1^{(r)}-\dot{\bar z}_1
  \in
  \{0.5143149884,\ -2,\ 7.4856850116\}.
  \label{eq:bf-hd-sgqf-example-Deltadot}
\end{equation}
The derived moment sensitivities are
\begin{align}
  \dot{\bar z}_1&=2.5=\frac52,
  &\dot v_1&=-2.5=-\frac52,
  \label{eq:bf-hd-sgqf-example-zbar-v-dot}\\
  \dot S_1&=14.5=\frac{29}{2},
  &\dot C_{xz,1}&=3.25=\frac{13}{4}.
  \label{eq:bf-hd-sgqf-example-S-Cxz-dot}
\end{align}
Let
\begin{equation}
  u_1=S_1^{-1}v_1=\frac{4}{43}\approx0.0930232558.
  \label{eq:bf-hd-sgqf-example-u}
\end{equation}
Then the one-step score ingredients are
\begin{align}
  \tr(S_1^{-1}\dot S_1)&=\frac{116}{43}\approx2.6976744186,
  \label{eq:bf-hd-sgqf-example-score-term1}\\
  2\dot v_1u_1&=-\frac{20}{43}\approx-0.4651162791,
  \label{eq:bf-hd-sgqf-example-score-term2}\\
  -u_1^2\dot S_1&=-\frac{232}{1849}\approx-0.1254732288.
  \label{eq:bf-hd-sgqf-example-score-term3}
\end{align}
The same-branch one-step score is therefore
\begin{equation}
  \dot{\widehat\ell}_1^{\rm FSGQ}
  =
  -\frac12\left[
    \frac{\dot S_1}{S_1}
    +2\dot v_1u_1
    -u_1^2\dot S_1
  \right]
  =-\frac{1948}{1849}
  \approx -1.0535424554.
  \label{eq:bf-hd-sgqf-example-ell-dot}
\end{equation}
For later propagation,
\begin{equation}
  \dot K_1
  =
  \dot C_{xz,1}S_1^{-1}-K_1\dot S_1S_1^{-1}
  =-\frac{42}{1849}
  \approx -0.0227149811,
  \label{eq:bf-hd-sgqf-example-Kdot}
\end{equation}
so
\begin{align}
  \dot m_1&=\dot m_1^-+\dot K_1v_1+K_1\dot v_1
  =-\frac{343}{3698}
  \approx -0.0927528394,
  \label{eq:bf-hd-sgqf-example-mdot}\\
  \dot P_1&=\dot P_1^- -\dot K_1S_1K_1-K_1\dot S_1K_1-K_1S_1\dot K_1
  =\frac{2353}{1849}
  \approx 1.2725797729.
  \label{eq:bf-hd-sgqf-example-Pdot}
\end{align}
This makes the chapter's limitation concrete.  Every quantity above is
deterministic and internally checkable on the declared branch, yet the nonlinear
observation law can still generate non-Gaussian posterior structure.  FixedSGQF
therefore gives a transparent Gaussian-surrogate value-and-gradient oracle, not a
full non-Gaussian posterior representation.

\paragraph{Implementation checklist for this oracle.}
A first implementation should be able to reproduce, up to declared floating-point
tolerance, at least
\(m_1^-\), \(P_1^-\), \(\bar z_1\), \(S_1\), \(C_{xz,1}\), \(K_1\),
\(m_1\), \(P_1\), \(\widehat\ell_1^{\rm FSGQ}\), and
\(\dot{\widehat\ell}_1^{\rm FSGQ}\).  This is the compact end-to-end numerical
oracle for both the value path and the derivative path.

This example is not a global correctness theorem.  Its role is narrower and
practical: a first BayesFilter implementation should be able to reproduce these
numbers before any same-branch derivative claim is trusted.

\section{General Analytical Derivative Framework For FixedSGQF}
\label{sec:bf-hd-sgqf-general-derivative}

The worked oracle fixes one concrete same-branch derivative instance.  The next
task is to state the general derivative framework so that the SGQF lane does not
depend on the toy example alone.  Differentiate with respect to one parameter
component \(\theta_i\), and let a dot denote \(\partial_i\).

\paragraph{Why the derivative is operationally delicate.}
The one-step scalar depends on the previous carried Gaussian, the covariance
factor used to place the cloud, the nonlinear transition and observation images,
the innovation covariance, and the posterior update.  The derivative must
therefore propagate through the entire recursion, and it remains meaningful only
when the value and derivative paths stay on the same declared branch.

\paragraph{Current-score objects versus propagation objects.}
For the current likelihood contribution, the derivative closes once
\(\dot v_t\) and \(\dot S_t\) are known.  The derivative
\(\dot C_{xz,t}\) is not needed to close the current score, but it is needed to
form \(\dot K_t\) and hence to propagate \(\dot m_t\) and \(\dot P_t\) to the
next time step.  A correct implementation must keep these two roles separate.

\subsection{Derivative Ledger}
\label{sec:bf-hd-sgqf-derivative-ledger}

The chain is
\begin{equation}
\begin{aligned}
  \theta
  &\longmapsto
  (m_{t-1},P_{t-1})
  \longmapsto
  C_{t-1}
  \longmapsto
  \chi_{t-1}^{(r)}
  \longmapsto
  a_t^{(r)}
  \longmapsto
  (m_t^-,P_t^-)
  \longmapsto
  C_t^- \\
  &\longmapsto
  \chi_t^{(r)}
  \longmapsto
  z_t^{(r)}
  \longmapsto
  (\bar z_t,S_t,C_{xz,t})
  \longmapsto
  (v_t,K_t,\widehat\ell_t^{\rm FSGQ},m_t,P_t).
\end{aligned}
  \label{eq:bf-hd-sgqf-derivative-chain}
\end{equation}
The sparse-grid nodes and weights do not move with \(\theta\).  They are saved
branch objects.  The Gaussian factors \(C_{t-1}\) and \(C_t^-\), however, do move
with \(\theta\), because they are functions of \(P_{t-1}\) and \(P_t^-\).

\subsection{Square-Root Branch}
\label{sec:bf-hd-sgqf-square-root-branch}

For the declared lower-triangular Cholesky branch \(P=LL^\top\), set
\begin{equation}
  A=L^{-1}\dot P L^{-\top}.
  \label{eq:bf-hd-sgqf-chol-A}
\end{equation}
Define the lower-triangular operator
\begin{equation}
  \Phi(A)_{jk}
  =
  \begin{cases}
  A_{jk}, & j>k,\\
  A_{jj}/2, & j=k,\\
  0, & j<k.
  \end{cases}
  \label{eq:bf-hd-sgqf-chol-Phi}
\end{equation}
Then
\begin{equation}
  \dot L=L\Phi(A).
  \label{eq:bf-hd-sgqf-chol-derivative}
\end{equation}
Every factor derivative in this SGQF lane uses this declared branch recipe.

\subsection{Prediction Sensitivities}
\label{sec:bf-hd-sgqf-prediction-sensitivities}

From the placed-point equation,
\begin{equation}
  \dot\chi_{t-1}^{(r)}=
  \dot m_{t-1}+\dot C_{t-1}\xi^{(r)}.
  \label{eq:bf-hd-sgqf-dot-pred-place}
\end{equation}
From the transition map,
\begin{equation}
  \dot a_t^{(r)}=
  D_xf_\theta(\chi_{t-1}^{(r)})\dot\chi_{t-1}^{(r)}
  +\partial_i f_\theta(\chi_{t-1}^{(r)}).
  \label{eq:bf-hd-sgqf-dot-transition}
\end{equation}
The predicted mean derivative is
\begin{equation}
  \dot m_t^-=
  \sum_{r=1}^M w_r\dot a_t^{(r)}.
  \label{eq:bf-hd-sgqf-dot-pred-mean}
\end{equation}
Let
\begin{equation}
  d_t^{(r)}=a_t^{(r)}-m_t^-,
  \qquad
  \dot d_t^{(r)}=\dot a_t^{(r)}-\dot m_t^-.
  \label{eq:bf-hd-sgqf-pred-centered}
\end{equation}
Then
\begin{equation}
  \dot P_t^-=
  \dot Q_\theta
  +
  \sum_{r=1}^M w_r
  \left\{\dot d_t^{(r)}d_t^{(r)\top}+d_t^{(r)}\dot d_t^{(r)\top}\right\}.
  \label{eq:bf-hd-sgqf-dot-pred-cov}
\end{equation}

\subsection{Observation Sensitivities And Innovation Score}
\label{sec:bf-hd-sgqf-observation-sensitivities}

On the predictive factor branch,
\begin{equation}
  \dot\chi_t^{(r)}=
  \dot m_t^-+\dot C_t^-\xi^{(r)}.
  \label{eq:bf-hd-sgqf-dot-obs-place}
\end{equation}
Observation sensitivities satisfy
\begin{equation}
  \dot z_t^{(r)}=
  D_xh_\theta(\chi_t^{(r)})\dot\chi_t^{(r)}
  +\partial_i h_\theta(\chi_t^{(r)}).
  \label{eq:bf-hd-sgqf-dot-obs-values}
\end{equation}
Then
\begin{align}
  \dot{\bar z}_t&=\sum_{r=1}^M w_r\dot z_t^{(r)},
  \label{eq:bf-hd-sgqf-dot-zbar}\\
  \dot v_t&=-\dot{\bar z}_t,
  \label{eq:bf-hd-sgqf-dot-v}\\
  \dot\Delta_t^{(r)}&=\dot z_t^{(r)}-\dot{\bar z}_t,
  \label{eq:bf-hd-sgqf-dot-Delta}\\
  \dot S_t&=\dot R_\theta+
  \sum_{r=1}^M w_r\left\{\dot\Delta_t^{(r)}\Delta_t^{(r)\top}+
  \Delta_t^{(r)}\dot\Delta_t^{(r)\top}\right\},
  \label{eq:bf-hd-sgqf-dot-S}\\
  \dot C_{xz,t}&=
  \sum_{r=1}^M w_r\left\{\dot\chi_t^{(r)}\Delta_t^{(r)\top}+
  (\chi_t^{(r)}-m_t^-)\dot\Delta_t^{(r)\top}\right\}.
  \label{eq:bf-hd-sgqf-dot-Cxz}
\end{align}
The one-step fixed-Gaussian innovation score is
\begin{equation}
  \partial_i\widehat\ell_t^{\rm FSGQ}
  =
  -\frac12\left[
    \tr(S_t^{-1}\dot S_t)
    +2\dot v_t^\top S_t^{-1}v_t
    -v_t^\top S_t^{-1}\dot S_tS_t^{-1}v_t
  \right].
  \label{eq:bf-hd-sgqf-score}
\end{equation}
This closes the current score role.  The propagation role then uses
\begin{equation}
  \dot K_t=
  \dot C_{xz,t}S_t^{-1}-K_t\dot S_tS_t^{-1}
  \label{eq:bf-hd-sgqf-dot-K}
\end{equation}
and
\begin{align}
  \dot m_t&=\dot m_t^-+\dot K_tv_t+K_t\dot v_t,
  \label{eq:bf-hd-sgqf-dot-m}\\
  \dot P_t&=\dot P_t^- -\dot K_tS_tK_t^\top-K_t\dot S_tK_t^\top-K_tS_t\dot K_t^\top.
  \label{eq:bf-hd-sgqf-dot-P}
\end{align}
These become the differentiated inputs to the next time step.

\section{Same-Branch Fixed-Cloud Scalar Contract}
\label{sec:bf-hd-sgqf-same-branch}

For the sparse-grid lane, the saved branch identity is not exhausted by the
symbol \(\mathcal C_{b,L}\).  A branch-valid finite-difference comparison must
reuse the full fixed ledger:
\begin{equation}
  \mathfrak B_{\rm FSGQ}=
  \{\mathcal C_{b,L},\ \text{one-dimensional rule family},\ \text{merge tolerance},\
  \text{zero-weight pruning rule},\ \text{node ordering},\ \text{factor branch},\
  \text{symmetrize-then-veto policy},\ \text{initial-condition policy}\}.
  \label{eq:bf-hd-sgqf-branch-record}
\end{equation}
For parameter component \(i\), the same-branch central difference is
\begin{equation}
  D_i(h)=
  \frac{\widehat\ell_T^{\rm FSGQ}(\theta+h e_i;\mathfrak B_{\rm FSGQ})-
        \widehat\ell_T^{\rm FSGQ}(\theta-h e_i;\mathfrak B_{\rm FSGQ})}{2h}.
  \label{eq:bf-hd-sgqf-central-diff}
\end{equation}
A row is branch-valid only if both perturbations reuse the same saved structural
branch record and realize the same accepted/failure stage-time pattern through
every likelihood contribution included in the comparison.  Equal cloud family but
different first failure time counts as a branch mismatch; equal first failure
time but different accepted-prefix pattern also counts as a branch mismatch.

That stricter logic matters because sparse-grid differentiation is not just
“differentiate some nearby Gaussian filter.”  It is “differentiate the same
fixed-cloud scalar.”  The later shared same-scalar chapter may import that
contract, but the sparse-grid lane must state it locally because the cloud,
merge policy, and branch-veto discipline are part of the lane-specific scalar
identity.

\paragraph{Frozen versus recomputed objects on a saved branch.}
The branch identity freezes structure, not numerical values.  Within a fixed
cloud branch,
\begin{equation}
\begin{aligned}
  \text{frozen structural objects:}&\quad
  \mathcal C_{b,L},\ \tau_{\rm merge},\ \tau_{\rm zero},\ \text{ordering},\ C(\cdot),\
  \text{preprocessing policy},\ \text{initial-condition policy},\ \epsilon_S,\epsilon_P,\\
  \text{parameter-dependent numerical objects:}&\quad
  m_t,P_t,C_t,m_t^-,P_t^-,C_t^-,\bar z_t,S_t,C_{xz,t},K_t,
  \widehat\ell_t^{\rm FSGQ}.
\end{aligned}
  \label{eq:bf-hd-sgqf-frozen-vs-variable}
\end{equation}
Thus same-scalar finite differences recompute the parameter-dependent numerical
objects by the same declared branch rule; they do not reuse old numerical values.
In particular, \(m_0(\theta\pm he_i)\) and \(P_0(\theta\pm he_i)\) are
recomputed by the same initial-condition policy.  What must remain unchanged is
the structural route by which those values are produced.

\paragraph{Finite-difference ladder and pass criterion.}
For each parameter component, use a ladder such as
\begin{equation}
  h\in\{10^{-1},10^{-2},10^{-3},10^{-4},10^{-5}\}\max\{1,|\theta_i|\}.
  \label{eq:bf-hd-sgqf-fd-ladder}
\end{equation}
For valid rows, compare the central difference against the analytical score
\(G_i\) by
\begin{equation}
  e_i(h)=
  \frac{|D_i(h)-G_i|}{\max\{1,|D_i(h)|,|G_i|\}}.
  \label{eq:bf-hd-sgqf-fd-relative-error}
\end{equation}
The diagnostic passes when at least two consecutive valid rows satisfy
\(e_i(h)\le 10^{-5}\) and the smaller-step row is no worse than ten times the
larger-step row.  Failure means one of three things is likely: the derivative
formula is wrong, the value and gradient paths are not using the same scalar, or
the declared branch is too nonsmooth or ill-conditioned at that point.

\paragraph{Implementation-facing FD checklist.}
A routine of the form
\texttt{same\_scalar\_fd\_check(}$i,h_{\rm ladder},\theta,\mathfrak B_{\rm FSGQ}$\texttt{)}
should:
\begin{enumerate}[leftmargin=0.28in]
  \item evaluate the value path at \(\theta\pm he_i\) using the same saved
  structural branch metadata,
  \item mark a row branch-invalid if either side changes the structural branch or
  the realized accepted/failure stage-time pattern,
  \item compute the central difference only for branch-valid rows,
  \item compare that central difference against the analytical score with
  \eqref{eq:bf-hd-sgqf-fd-relative-error}, and
  \item report whether two consecutive valid rows satisfy the declared tolerance
  rule.
\end{enumerate}
The equations above remain the authoritative mathematical definition; this list is
the operational review contract.

\section{Lane-Specific Validation Burden}
\label{sec:bf-hd-sgqf-validation}

Fixed SGQF should be validated as an approximate deterministic likelihood lane,
not as an exact posterior filter.  The minimum validation suite is:
\begin{enumerate}[leftmargin=0.28in]
  \item \textbf{Construction checks.}  Verify duplicate-node accumulation,
  weight totals, signed-weight counts, and deterministic cloud reproduction for
  small \((b,L)\) cases.
  \item \textbf{Linear-Gaussian recovery.}  For affine \(f_\theta,h_\theta\), the
  Gaussian projection recursion should match the Kalman filter up to
  floating-point tolerance.
  \item \textbf{Worked scalar oracle.}  Reproduce the one-step values in
  Section~\ref{sec:bf-hd-sgqf-worked-example} exactly enough that the cloud,
  factor, innovation, and likelihood assembly are demonstrably aligned.
  \item \textbf{Same-branch finite differences.}  Use
  \eqref{eq:bf-hd-sgqf-central-diff} on representative parameters and step
  sizes while logging branch-valid versus branch-invalid rows.
  \item \textbf{Signed-weight stability.}  Record the minimum eigenvalues of
  symmetrized \(P_t^-\), \(S_t\), and \(P_t\), plus any branch veto.
  \item \textbf{Memory and runtime scaling.}  Record \(M_{b,L}\), model
  evaluations, wall time, and peak memory as \(b\), \(L\), and \(T\) change.
  \item \textbf{Accuracy ladder.}  Compare fixed SGQF against tensor-product GHQ
  in small dimension and against other controlled references where available.
\end{enumerate}

A useful test-model ladder mirrors the p47 source burden:
\begin{itemize}[leftmargin=0.28in]
  \item \textbf{Model A:} a linear-Gaussian recovery case,
  \item \textbf{Model B:} the scalar quadratic-observation cell from
  Chapter~\ref{ch:bf-highdim-nonlinear-foundations},
  \item \textbf{Model C:} a cloud-sensitive two-dimensional nonlinear observation
  where duplicate merging and fourth moments matter,
  \item \textbf{Model D:} a small high-dimensional block-local model where signed
  weights, active-dimension growth, and covariance checks become visible,
  \item \textbf{Model E:} a deliberately hard all-coordinate interaction model to
  expose when sparse-grid savings stop compensating for Gaussian-surrogate bias,
  \item \textbf{Model F:} a Jia--Xin--Cheng-style orbit / sparse-grid benchmark
  that distinguishes sparse-grid exactness claims from posterior-accuracy claims.
\end{itemize}
The first three models should be understood concretely, not only as names:
\begin{align}
  \text{Model A: }&
  X_t=A_\theta X_{t-1}+\eta_t,
  \quad
  \eta_t\sim\calN(0,Q_\theta),
  \quad
  Y_t=H_\theta X_t+\varepsilon_t,
  \quad
  \varepsilon_t\sim\calN(0,R_\theta),
  \label{eq:bf-hd-sgqf-model-A}\\
  \text{Model B: }&
  X_t=\rho X_{t-1}+\eta_t,
  \quad
  \eta_t\sim\calN(0,q),
  \quad
  Y_t=X_t^2+\varepsilon_t,
  \quad
  \varepsilon_t\sim\calN(0,r),
  \label{eq:bf-hd-sgqf-model-B}\\
  \text{Model C: }&
  X_t=\rho X_{t-1}+\eta_t,
  \quad
  \eta_t\sim\calN(0,qI_2),
  \quad
  Y_t=\theta^2(X_{t,1}^2+X_{t,2}^2)+\varepsilon_t,
  \quad
  \varepsilon_t\sim\calN(0,r).
  \label{eq:bf-hd-sgqf-model-C}
\end{align}
Model A checks exact-reference recovery, Model B exposes posterior-shape failure
of one-Gaussian surrogates, and Model C catches duplicate-node or signed-weight
assembly errors because the relevant fourth moments change when the cloud is
built incorrectly.

The implementation report for this lane should log at least:
\begin{equation}
  \mathcal R_{\rm SGQF}=
  \{M_{b,L},\ \text{signed-weight count},\ \min\lambda(P_t^-),\ \min\lambda(S_t),\
    \min\lambda(P_t),\ \text{FD table},\ \text{branch-veto count},\
    \text{wall time},\ \text{peak memory}\}.
  \label{eq:bf-hd-sgqf-report-object}
\end{equation}
Accuracy claims should be made only against a named comparator:
\begin{equation}
  \text{error}
  =
  \begin{cases}
  |\widehat\ell_T-\ell_T^{\rm Kalman}|, & \text{linear-Gaussian test},\\
  |\widehat\ell_T-\ell_T^{\rm dense}|, & \text{small dense-quadrature test},\\
  \text{RMSE or posterior diagnostic against PF/TT reference}, & \text{nonlinear large test}.
  \end{cases}
  \label{eq:bf-hd-sgqf-comparator-errors}
\end{equation}
These are not mere software counters.  They record the places where the declared
sparse-grid scalar can fail mathematically or operationally: duplicate merging,
signed-weight stability, covariance branch validity, same-branch derivative
parity, and cost growth.

Passing this suite supports the claim that the implementation evaluates and
differentiates the declared fixed-cloud surrogate scalar.  It does not prove
exact posterior accuracy, HMC convergence, or default readiness.

\section{Implementation Contract For FixedSGQF}
\label{sec:bf-hd-sgqf-implementation-contract}

A compact implementation contract makes the SGQF lane auditable without forcing
the reader back to the source note.  The key runtime objects are:
\begin{center}
\small
\begin{tabular}{p{0.18\linewidth}p{0.30\linewidth}p{0.16\linewidth}p{0.24\linewidth}}
\toprule
object & definition & shape & reuse \\
\midrule
\(\xi^{(r)}\) & standardized sparse-grid node & \(n_x\) & value and gradient \\
\(w_r\) & accumulated signed weight after duplicate merge & scalar & value and gradient \\
\(C_{t-1},C_t^-\) & covariance factors on declared branch & \(n_x\times n_x\) & value, gradient, diagnostics \\
\(\chi_{t-1}^{(r)}\) & previous placed point & \(n_x\) & transition value and derivative \\
\(a_t^{(r)}\) & transition image & \(n_x\) & prediction moments \\
\(\chi_t^{(r)}\) & predictive placed point & \(n_x\) & observation value and derivative \\
\(z_t^{(r)}\) & observation image & \(n_y\) & observation moments \\
\(\bar z_t,S_t,C_{xz,t}\) & Gaussian projection moments & \(n_y\), \(n_y\times n_y\), \(n_x\times n_y\) & likelihood and update \\
\(v_t,K_t,u_t\) & innovation, gain, solve variable with \(S_tu_t=v_t\) & \(n_y\), \(n_x\times n_y\), \(n_y\) & likelihood, update, gradient \\
\(\dot\chi,\dot a,\dot z\) & point sensitivities & matching value objects & gradient only \\
\(\dot m,\dot P,\dot S,\dot C_{xz}\) & moment and filter sensitivities & matching value objects & gradient and next step \\
data record & \(y_{1:T}\), preprocessing, missing-data policy if any & metadata & value, gradient, finite difference \\
initial law & \(m_0(\theta),P_0(\theta),\dot m_0,\dot P_0\) & \(n_x,n_x\times n_x\) & value and gradient start \\
branch record & cloud, rules, merge tolerance, zero-weight rule, node ordering, preprocessing policy, initial-condition policy, factor branch, symmetrize-then-veto rule, thresholds, accepted/failure stage-time pattern & metadata & same-scalar finite difference \\
\bottomrule
\end{tabular}
\end{center}

\subsection{Input, Output, Invariant, And Failure Contract}
\label{sec:bf-hd-sgqf-io-contract}

The full lane-level routine has inputs
\begin{equation}
\begin{aligned}
  \mathcal I_{\rm in}=
  (&y_{1:T},\theta,L,\{I_\ell\}_{\ell=1}^{L+b-1},
  m_0(\theta),P_0(\theta),\\
  &\dot m_0^{(1:p)}(\theta),\dot P_0^{(1:p)}(\theta),
  f_\theta,h_\theta,Q_\theta,R_\theta,\\
  &D_xf_\theta,\partial_{1:p}f_\theta,
  D_xh_\theta,\partial_{1:p}h_\theta,
  \dot Q_\theta^{(1:p)},\dot R_\theta^{(1:p)},\\
  &\text{merge tolerance},\text{factor map},\text{branch thresholds}).
\end{aligned}
  \label{eq:bf-hd-sgqf-input-contract}
\end{equation}
It returns either
\begin{equation}
  \left(
  \widehat\ell_T^{\rm FSGQ}(\theta),
  \nabla_\theta\widehat\ell_T^{\rm FSGQ}(\theta),
  \{m_t,P_t\}_{t=0}^T,
  \mathcal D_{\rm diag}
  \right)
  \label{eq:bf-hd-sgqf-output-success}
\end{equation}
or a failure record
\begin{equation}
  \mathcal F=(\text{time},\text{stage},\text{reason},\text{diagnostic values}).
  \label{eq:bf-hd-sgqf-output-failure}
\end{equation}
The invariants are:
\begin{align}
  \sum_{r=1}^M w_r&=1
  &&\text{up to construction tolerance},
  \label{eq:bf-hd-sgqf-invariant-weight-sum}\\
  P_t^-&=(P_t^-)^\top,\quad
  S_t=S_t^\top,\quad
  P_t=P_t^\top
  &&\text{after explicit symmetrization},
  \label{eq:bf-hd-sgqf-invariant-symmetry}\\
  S_t&\succ0
  &&\text{on every accepted likelihood step},
  \label{eq:bf-hd-sgqf-invariant-positive-S}\\
  \mathcal B_{\rm value}&=\mathcal B_{\rm gradient}
  &&\text{same saved branch}.
  \label{eq:bf-hd-sgqf-invariant-same-branch}
\end{align}
The routine fails rather than repairs if
\begin{equation}
  \lambda_{\min}(S_t)<\epsilon_S,\qquad
  \lambda_{\min}(P_t^-)<\epsilon_P,\qquad
  \lambda_{\min}(P_t)<\epsilon_P,
  \label{eq:bf-hd-sgqf-failure-eigen}
\end{equation}
or if the declared factor map cannot be evaluated on the symmetrized covariance.
A smoothed repair would define a different scalar and therefore would need a new
derivative contract.

\subsection{One-Step Value-Path Contract}
\label{sec:bf-hd-sgqf-one-step-value-contract}

At one accepted time step, the value path is the ordered map
\begin{equation}
\begin{aligned}
  (m_{t-1},P_{t-1},y_t,\theta,\mathcal C)
  &\longmapsto C_{t-1}
  \longmapsto \chi_{t-1}^{(r)}
  \longmapsto a_t^{(r)}
  \longmapsto (m_t^-,P_t^-) \\
  &\longmapsto C_t^-
  \longmapsto \chi_t^{(r)}
  \longmapsto z_t^{(r)}
  \longmapsto (\bar z_t,S_t,C_{xz,t}) \\
  &\longmapsto (v_t,\widehat\ell_t^{\rm FSGQ},K_t)
  \longmapsto (m_t,P_t).
\end{aligned}
  \label{eq:bf-hd-sgqf-one-step-value-map}
\end{equation}
The accepted map is conditioned on two successful branch checks: the predictive
covariance \(P_t^-\) must pass the declared factor and positive-definite test
before observation-point placement, and the innovation covariance \(S_t\) must
pass the declared branch test before the Gaussian projection update is accepted.
The carried outputs propagated to the next step are only \((m_t,P_t)\) together
with the global fixed branch record \(\mathcal B\).  Point clouds, moment arrays,
gains, and diagnostics may be cached for derivative reuse, but they are not part
of the carried filtering state itself.

\subsection{Default Numerical Constants}
\label{sec:bf-hd-sgqf-defaults}

The constants below are defaults for diagnostics, not claims of universal
validity:
\begin{center}
\small
\begin{tabular}{p{0.24\linewidth}p{0.22\linewidth}p{0.44\linewidth}}
\toprule
quantity & default & role \\
\midrule
duplicate merge tolerance & \(10^{-12}\max(1,\|\xi\|_\infty)\) & groups standardized nodes that should be identical under the declared rule \\
weight zero tolerance & \(10^{-14}\) & removes accumulated numerical-zero weights after merge \\
\(\epsilon_S\) & \(10^{-10}\) & veto threshold for innovation covariance positive definiteness \\
\(\epsilon_P\) & \(10^{-10}\) & veto threshold for carried and predictive covariance factors \\
\(\epsilon_\ell\) & \(10^{-3}\max(1,|\widehat\ell_T|)\) & pilot level-ladder change threshold \\
finite-difference relative tolerance & \(10^{-5}\) & same-scalar gradient diagnostic \\
maximum branch-invalid FD rows & 0 in final audit & branch changes mean the derivative is not being tested \\
\bottomrule
\end{tabular}
\end{center}
These values are deliberately conservative.  They are meant to reveal when the
fixed branch is fragile.  They are not tuning knobs for making a failing branch
look smooth.

\section{End-To-End Mathematical Algorithm}
\label{sec:bf-hd-sgqf-end-to-end-algorithm}

\fbox{\begin{minipage}{0.94\linewidth}
\textbf{FixedSGQF value and gradient, with all branch exits.}

\begin{enumerate}[leftmargin=0.22in]
\item Build the sparse-grid cloud from \(b=n_x\), level \(L\), and
\(\{I_\ell\}\).  Save \(\mathcal C=\{(\xi^{(r)},w_r)\}_{r=1}^M\).
\item Check the weight-sum invariant.  If it fails, return
\(\mathcal F=(0,\text{cloud},\text{weight-sum failure},\cdots)\).
\item Initialize \(\widehat\ell_T^{\rm FSGQ}\gets0\),
\(G_i\gets0\) for \(i=1,\ldots,p\), and compute
\(m_0,P_0,\dot m_0^{(i)},\dot P_0^{(i)}\).
\item For \(t=1,\ldots,T\), symmetrize \(P_{t-1}\).  If the factor map
fails, return a factor-branch failure.  Otherwise set
\(C_{t-1}=C(P_{t-1})\) and compute the matching factor derivatives.
\item Place previous points, compute transition images, and form
\((m_t^-,P_t^-)\) together with their derivatives.
\item Symmetrize \(P_t^-\).  If the predictive covariance fails the branch,
return predictive covariance failure.
\item Factor \(P_t^-\), place observation points, evaluate observation values,
and form \((\bar z_t,S_t,C_{xz,t},v_t)\).
\item Symmetrize \(S_t\).  If the innovation covariance fails the branch,
return innovation covariance failure.
\item Only on that accepted innovation branch, form
\(\widehat\ell_t^{\rm FSGQ}\), update
\(\widehat\ell_T^{\rm FSGQ}\), and accumulate the componentwise score.
\item Form \(K_t\), then propagate \(m_t,P_t\) and their sensitivities.
If the carried covariance branch fails, return carried covariance failure.
\item Return
\((\widehat\ell_T^{\rm FSGQ},G_{1:p},\{m_t,P_t\},\mathcal D_{\rm diag})\).
\end{enumerate}
\end{minipage}}

The algorithm is deliberately mathematical rather than Pythonic.  Its purpose is
to say what an implementation must compute, save, and refuse to change if the
reported gradient is to be the derivative of the reported value.

\section{Exports To Later Chapters}
\label{sec:bf-hd-sgqf-validation-exports}

This SGQF continuation chapter exports five specific objects forward:
\begin{enumerate}[leftmargin=0.28in]
  \item the fixed-cloud value-path recursion,
  \item the one-step numeric oracle,
  \item the lane-specific branch-identity record,
  \item the branch-valid finite-difference contract for SGQF,
  \item the SGQF validation and non-claim boundaries.
\end{enumerate}
The later fixed-branch chapter may now treat the sparse-grid lane as an already
specified scalar object rather than re-explaining the cloud mechanics.  The
closing validation/promotion chapter may consume the SGQF diagnostics and veto
logic without needing to rederive the lane’s runtime map.
